\subsection{Workflow guidance, scientific judgment, and outlook}

The primary comparison shows that Cloudra can improve the same backbone's research delivery by keeping chemical objects, computational results, and scientific requirements connected throughout exploration. The consistent Hard-task improvement and the Normal-task gains place this contribution in comparative stereochemical studies and the completion of multistep research. The complete 50-task decomposition locates gains in both evidence closure and target accuracy; in Hard tasks, their weighted contributions are +0.1353 and +0.1274. HH1 and HI1 additionally show why evidence closure matters even after a reference-consistent barrier has been obtained: the numerical result and the path supporting its scientific attribution are separate achievements.

The archived cases suggest three requirements for an autonomous research process. First, the identity of a scientific object must persist beyond the current language description. A reactant species can contain several conformers, and a requested reaction step can produce an optimized saddle belonging to a different route. HA8 requires the energy denominator to remain associated with the stable species while the IRC retains its actual local endpoints. NU6 requires a newly optimized stationary point to be assigned from its resulting connectivity rather than from the intended name of the initial guess. These are global relations among results accumulated at different times, consistent with the role of external planning and verification feedback in language-model workflows. A persistent research graph gives the model a place to inspect and update them as each calculation changes what is known.

Second, feedback must identify the level at which the research should change. A missing endpoint optimization calls for a different action from a saddle that repeatedly relaxes to another route, and neither is equivalent to a method that cannot describe the target electronic state. Language models generate plausible next actions from their current context; a chemically familiar description can remain coherent even when the numerical or path evidence has changed. External measurements and structured evidence can expose that mismatch, but the agent must still decide whether to repair a calculation, revise the pathway hypothesis, change a reference assignment, or withhold the requested conclusion. The Record comparison shows that making tools available is insufficient to ensure that these relationships remain active in the final argument. The Brief-errors comparison further shows that the useful content of feedback depends on the decision being made, complementing evidence on the importance of external feedback for reasoning correction.

Third, completing a local proof and recognizing the scope of a scientific claim need to be supported together. Brief-errors' higher Impossible mean came from particular target judgments in IA1 and IA9, with the reverse result in IG8. Its IC3 graph provides a more detailed behavioral example: the run retained a characterized closed-shell branch while refusing to treat it as the requested lowest open-shell singlet or to report its relative energy as the target gap. This distinction is a valuable form of scientific judgment. It preserves what exploration established without allowing a convenient result to replace a different research question. The crossed outcomes motivate a task-dependent hypothesis: detailed diagnostics can help complete a valid local proof, while continuing to repair that proof may be unhelpful when the method or target itself needs reconsideration. The present final-result comparisons do not establish the causal effect of diagnostic length or measure the internal cognitive mechanism responsible.

The score and claim records also reveal where the evaluator's scientific identities need refinement. Record/IA5 correctly rejected the requested disrotatory path, but the archived adjudication records that a matcher selected another positive path with the same endpoints. Brief-errors/IC3 received accuracy 0 despite explicitly distinguishing the requested open-shell state from its separate closed-shell branch. We retain the released scores, while reading them together with the saved claims. These cases indicate that molecular connectivity alone cannot identify a reaction channel or an electronic-state target. The evaluation needs to retain the same state and pathway distinctions that the research workflow is intended to preserve.

P2 contains one retained run per task and configuration; task-cluster bootstrap intervals describe variation in the task mix rather than independent repeat-run variability. The complete organized table covers all 200 runs and their three score components. Some Hard-task accuracy targets include qualitative channel identity, so the reported gain concerns research completion quality as well as numerical barriers. Repeated evaluation on the same tasks would further characterize run-to-run variability and the conditions under which a particular research route is selected.

The Impossible tier is defined by the allowed-method contract, rather than a universal claim that the corresponding chemistry cannot occur. Its component values also use negative-target rules: self-consistency 1 on an explicit negative conclusion is distinct from a positive TS--IRC--endpoint chain. The 14 Impossible cases with self-consistency 1 and accuracy 0 should consequently be interpreted through their target claims and provenance. NU6's matching energy fingerprint identifies a related problem for accessible pathways: an energy difference can acquire different scientific meanings when its saddle identity and the requested step are not maintained together. These cases motivate persistent checks of what each calculation actually answers as the research hypothesis changes.

Timing and effort are likewise reported at their observed units. The 12 pilot batches measure real elapsed time for four tasks at concurrency two, including the original failures and stopping outcomes. The P2 effort archive covers selected deep-dive studies, with 15 tasks shared across configurations; cumulative job-hours represent recorded quantum-chemical process durations, with a documented native-log reconstruction for Standard/M1. They are distinct from core-hours, financial cost, and end-to-end P2 research latency. Concurrent calculations shared computational resources, and CPU queueing may have delayed individual results. The measurements therefore support a comparison of the observed workflows and investments, while further separating queueing, calculation, and agent-active time would clarify where elapsed time is spent.

The delivery cases identify a concrete systems direction alongside the scientific one. Waiting for a calculation, receiving its completion event, resuming the research, and assembling the final artifact must share a continuous task lifecycle. Standard's missing graphs followed premature completion while calculations were still pending; Guided/HI1 also exposed a notification that did not wake a waiting session. Preserving an unfinished goal across these transitions allows newly completed computations to become usable research information. At the scientific level, future work should make repeated returns to the same stationary-point region visible and help the agent decide when another local attempt is unlikely to resolve the target pathway. Better candidate selection and explicit revision of research hypotheses can then complement the improved preservation of evidence.

Cloudra alleviates the difficulties of direct geometric reasoning and manipulation by language models, and its conformational sampling has improved substantially over earlier versions. Highly flexible systems nevertheless remain challenging. As conformational diversity grows, identifying relevant reactant arrangements, intermediates, and pathways can require substantially more sampling and quantum-chemical evaluation. More efficient conformational exploration and better allocation of calculations among competing candidates are therefore important directions for further development.

Our longer-term goal extends beyond small-molecule mechanism calculations to a broader range of computational chemistry. Support for multiple computational engines, shared execution infrastructure, and persistent scientific records already provides a basis for that expansion. Extending coverage will require appropriate representations, methods, and validation criteria for each new class of problem. We aim for Cloudra to become a powerful research tool for computational chemists and a computational specialist for experimental teams. Cloudra will be released as open-source software under the GNU Affero General Public License (AGPL), supporting inspection, reuse, and community development of the system.
